\title{\abstraction{}: Enhancing Observability \& Programmability in Smart Spaces}

\author{
{\rm Anna Karanika \hspace{1em} Kai-Siang Wang \hspace{1em} Han-Ting Liang \hspace{1em} Shalni Sundram \hspace{1em} Indranil Gupta}
\\\\
University of Illinois Urbana-Champaign
} % end author

\maketitle


\begin{abstract}

While RPCs form the bedrock of systems stacks, we posit that IoT device collections in smart spaces like homes, warehouses, and office buildings---which are all ``user-facing''---require a more expressive abstraction.
Orthogonal to prior work, which improved the {reliability} of IoT communication, our work focuses on improving the {\it observability} and {\it programmability} of IoT actions.
We present the \abstraction{} (Request-Acknowledge-Start-Complete) abstraction, which provides acknowledgments at critical points after an IoT device action is initiated. \abstraction{} is a better fit for
IoT actions, which naturally vary in length {\it spatially} (across devices) and {\it temporally} (across time, for a given device).
\abs{} also enables the design of several new features: predicting action completion times accurately, detecting failures of actions faster, allowing fine-grained dependencies in programming, and scheduling.
\abs{} is intended to be implemented atop today's available RPC mechanisms, rather than as a replacement. We integrated \abs{} into a popular and open-source IoT framework called Home Assistant.
Our trace-driven evaluation finds that \abstraction{} meets latency SLOs,
especially for long actions that last  O(mins), which are common in smart spaces.
Our scheduling policies for home automations (e.g., routines) outperform state-of-the-art counterparts by $10\%$-$55\%$.

\end{abstract}
